% The preamble every paper of the lab shares: \input{../../../papers/preamble}
% from models/<model>/paper/, \input{../preamble} from papers/<name>/.
% Boards come from chessboard and xskak, plots from pgfplots. A paper prints no
% number by hand: `lab paper` writes every number, table and figure under the
% paper's generated/ folder, and paper.tex says what they mean.
\usepackage[LSBC4,T1]{fontenc}
\usepackage[utf8]{inputenc}
\usepackage{lmodern}
\usepackage[margin=1in]{geometry}
\usepackage{microtype}
\usepackage{amsmath}
\usepackage{booktabs}
\usepackage{tabularx}
\usepackage{array}
\usepackage{xcolor}
\usepackage[obeyspaces,spaces]{url}
\usepackage{xskak}
\usepackage{pgfplots}
\usepgfplotslibrary{groupplots}
\usetikzlibrary{positioning,arrows.meta,calc,fit}
\pgfplotsset{compat=1.18}
\usepackage{caption}
\usepackage{subcaption}
\usepackage{enumitem}
\usepackage[section]{placeins}
\usepackage[hidelinks]{hyperref}

% A generated value by its key: generated/values.tex sets them, \V{key} prints
% one, and a key nobody set stops the build.
\makeatletter
\newcommand{\setvalue}[2]{\expandafter\gdef\csname labvalue@#1\endcsname{#2}}
\newcommand{\V}[1]{\@ifundefined{labvalue@#1}%
  {\PackageError{lab}{No generated value named '#1'}{Run `lab paper` for this model again.}}%
  {\csname labvalue@#1\endcsname}}
\makeatother

% Chart colours. A run keeps its colour in every figure of every paper.
\definecolor{labBlue}{HTML}{2A78D6}
\definecolor{labOrange}{HTML}{EB6834}
\definecolor{labAqua}{HTML}{1BAF7A}
\definecolor{labYellow}{HTML}{EDA100}
\definecolor{labMagenta}{HTML}{E87BA4}
\definecolor{labGreen}{HTML}{008300}
\definecolor{labViolet}{HTML}{6250D6}
\definecolor{labGray}{HTML}{898781}
\definecolor{labInk}{HTML}{0B0B0B}
\definecolor{labSecondary}{HTML}{52514E}
\definecolor{labGrid}{HTML}{E1E0D9}
\definecolor{labAxis}{HTML}{C3C2B7}
\definecolor{labWash}{HTML}{CDE2FB}

\pgfplotsset{
  lab/.style={
    font=\sffamily\footnotesize,
    axis line style={labAxis},
    tick style={labAxis},
    tick label style={labSecondary},
    label style={labSecondary},
    title style={labInk, font=\sffamily\small\bfseries, align=left},
    grid style={labGrid, line width=0.3pt},
    major grid style={labGrid, line width=0.3pt},
    axis lines*=left,
    tick align=outside,
    clip=false,
    legend style={draw=none, fill=none, font=\sffamily\footnotesize, text=labSecondary,
      cells={anchor=west}, legend columns=-1, column sep=6pt, at={(0,1.03)}, anchor=south west},
    every axis plot/.append style={line width=1.1pt, line join=round, line cap=round},
    scaled ticks=false,
    /pgf/number format/1000 sep={,},
  },
  lab marks/.style={mark=*, mark size=1.9pt, mark options={draw=white, line width=0.6pt}},
}

\setchessboard{boardfontencoding=LSBC4, boardfontfamily=skaknew, blackfieldcolor=black!18,
  whitefieldcolor=white, setfontcolors, showmover=false, label=true, labelfontsize=6pt,
  boardfontsize=12pt}

% A move or line in UCI, a position in FEN (breakable at a slash), and a run's name.
\newcommand{\uci}[1]{\texttt{#1}}
\newcommand{\fen}[1]{{\footnotesize\path{#1}}}
\newcommand{\run}[1]{\texttt{#1}}
\newcolumntype{L}{>{\raggedright\arraybackslash}X}
\newcolumntype{R}{>{\raggedleft\arraybackslash}X}

\setlength{\parskip}{0.4em}
\setlength{\parindent}{0pt}
\captionsetup{font=small, labelfont=bf}
